\chapter{Related Work}
\label{chap:related-work}

\section{Travel and Conversational Recommendation}

Travel recommendation differs from recommendation in domains where a user
selects one independent item. A travel decision may depend on location,
duration, budget, mobility, season, group composition, interests, and relations
between several points of interest. It is also common for users to refine these
requirements through conversation rather than state them completely at the
beginning.

Conversational recommender systems support preference elicitation, questions,
feedback, and explanation through multi-turn interaction. Jannach et
al.~\cite{jannach2021conversational} organize conversational recommendation
according to supported intents, background knowledge, interaction strategy,
and evaluation. Gao et al.~\cite{gao2021conversational} further identify
preference elicitation, multi-turn strategy, dialogue understanding,
exploration, and evaluation as central challenges. These observations are
directly relevant to travel assistance, where a user may change constraints or
refer to earlier recommendations.

Traditional conversational recommenders frequently optimize recommendation
accuracy over a predefined item catalog. The present work instead treats the
assistant as both a recommender and an evidence-grounded information system.
The goal is not only to rank a hotel or attraction, but also to explain the
choice using retrieved tourism text, preserve conversation state, and disclose
when requested evidence is absent.

\section{Retrieval-Augmented Generation}

Lewis et al.~\cite{lewis2020rag} introduced a general RAG formulation that
combines parametric generation with retrieved non-parametric memory for
knowledge-intensive NLP tasks. RAG makes it possible to update or inspect the
external knowledge source independently of the generator. This is valuable for
tourism, where facts are distributed across sources and users benefit from
source attribution.

Nevertheless, standard retrieve-then-generate pipelines are sensitive to
retrieval quality. Irrelevant context can distract the generator, while missing
context can lead to incomplete answers or unsupported use of parametric
knowledge. The proposed system therefore exposes intermediate retrieval stages,
checks requirement-level evidence coverage, and gates generation according to
answer readiness.

\section{Sparse and Dense Retrieval}

Sparse lexical retrieval remains effective for exact entity names, rare terms,
and queries whose wording overlaps with the corpus. BM25 is a widely used
probabilistic lexical ranking method and provides a strong first-stage retrieval
signal \cite{robertson2009bm25}. In tourism search, lexical matching is useful
for names such as attractions, markets, districts, and transport terminals.

Dense retrieval maps queries and passages into a continuous vector space.
Karpukhin et al.~\cite{karpukhin2020dpr} demonstrated that learned dense passage
representations can be effective for open-domain question answering. Sentence-
BERT provides semantically meaningful sentence embeddings that can be compared
efficiently using cosine similarity \cite{reimers2019sentencebert}. Dense
retrieval can therefore match related expressions even when they share few
words, such as ``family-friendly places'' and ``activities suitable for
children''.

Sparse and dense retrieval fail in different ways. Dense representations may
underweight an uncommon exact name, whereas BM25 may miss a semantic paraphrase.
This thesis combines both signals rather than selecting only one retriever.

\section{Rank Fusion and Neural Reranking}

Reciprocal Rank Fusion combines rankings using reciprocal rank contributions
without requiring their raw scores to be calibrated. Cormack et
al.~\cite{cormack2009rrf} showed that RRF can improve over individual rankers
and other rank-fusion approaches. Its simplicity is useful when BM25 scores and
vector similarities have different scales.

After first-stage retrieval, a neural reranker can examine query and passage
jointly. Nogueira and Cho~\cite{nogueira2019bert} demonstrated BERT-based
passage reranking for information retrieval. CrossEncoders are computationally
more expensive than independently encoded embeddings but can model direct token
interactions. The present system therefore applies a CrossEncoder only to a
bounded fused shortlist.

The main difference from a basic hybrid pipeline is that the proposed method
also uses controlled metadata signals, optional filter relaxation, comparison
balancing, evidence coverage, and recovery retrieval. These mechanisms address
travel-specific problems such as geographic ambiguity and multi-destination
requests.

\section{Corrective and Agentic Retrieval}

Corrective RAG methods recognize that generation quality depends on whether the
retrieved evidence is adequate. Yan et al.~\cite{yan2024crag} proposed a
retrieval evaluator that triggers different corrective actions when retrieved
documents are unreliable. The current thesis adopts the general principle of
evaluating retrieval before generation, but uses a bounded local recovery
process rather than unrestricted web search.

The planner first identifies evidence requirements and focused retrieval tasks.
After reranking, a checker labels each requirement as covered, partially
covered, or not covered. Up to three targeted recovery queries may then be
executed. This bounded design makes the process traceable and limits latency and
cost. The diagnostic layer also avoids claiming a corpus gap when a requirement
was never independently probed.

\section{Personalization and Memory}

Conversational personalization requires both long-term and short-term context.
A durable preference such as avoiding large tour groups may remain useful
across trips, whereas a five-day duration applies only to one conversation. If
these are stored together, old trip constraints can incorrectly influence later
recommendations.

The proposed system therefore maintains persistent user memory and temporary
conversation state as separate structures. Memory is used selectively during
ranking and answer generation; it does not automatically become a hard filter.
This design is consistent with conversational-recommendation research that
emphasizes dynamic preference elicitation and contextual interaction
\cite{jannach2021conversational,gao2021conversational}, while adding explicit
controls for memory lifetime and retrieval behavior.

\section{Evaluation of RAG Systems}

RAG evaluation is inherently multi-dimensional because the retriever and
generator can fail independently. RAGAS evaluates dimensions including
retrieved context and generation faithfulness without requiring complete
reference answers \cite{es2024ragas}. ARES similarly addresses automated RAG
evaluation with model-based judges \cite{saadfalcon2024ares}. These frameworks
motivate stage-level measurement, but a domain-specific system still requires
metrics aligned with its own query schema, personalization behavior, and
retrieval depth.

LLM-as-a-judge offers scalable evaluation of open-ended responses. Zheng et
al.~\cite{zheng2023llmjudge} report substantial agreement between strong LLM
judges and human preferences, while also identifying position, verbosity, and
self-enhancement biases. Consequently, this thesis uses a model separate from
the answer generator for initial judgments and retains human review rather than
treating judge scores as unquestionable ground truth.

\section{Research Gap}

Existing work provides strong foundations for conversational recommendation,
dense retrieval, rank fusion, neural reranking, corrective retrieval, and
automated evaluation. However, these elements do not by themselves define an
end-to-end personalized travel assistant that can explain stage-level failures.
The gap addressed in this thesis is their integration into a single traceable
pipeline with:

\begin{itemize}
    \item separate persistent memory and temporary trip state;
    \item structured intent, operation, constraint, and facet parsing;
    \item metadata-aware hybrid retrieval with safe geographic behavior;
    \item bounded requirement-level planning and recovery;
    \item explicit complete, partial, and insufficient readiness; and
    \item a unified evaluation trace from query understanding to cost and final
    answer quality.
\end{itemize}

